The standard model of particle physics is currently the most comprehensive and experimentally validated framework describing the fundamental particles and their interactions. It successfully unifies three of the four fundamental forces: electromagnetism; the weak nuclear force; and the strong nuclear force, into a single theoretical structure. Within this model, all known elementary particles are categorized into two primary classes: fermions and bosons. Fermions, which include quarks and leptons, are the building blocks of matter. Quarks combine to form mesons (e.g $J/\psi$) and baryons (e.g protons and neutrons), while leptons are elementary particles which include electrons, muons, tau and their associated flavor neutrinos. In contrast, bosons are force carriers that mediate interactions between fermions, transmitting the fundamental forces through the exchange of virtual particles. All the particles are listed in the illustration in Fig. \ref{fig:standard_model}.

\begin{figure}[h!] 
    \centering % Centers the image and caption
    \includegraphics[width=0.8\textwidth]{\main/illustrations/standard_model.png}
    \caption{The table describing the standard model. Quarks outlined in purple and leptons in green \cite{noauthor_standard_nodate}.}
    \label{fig:standard_model}
\end{figure}


Among the bosons, the photon serves as the mediator of the electromagnetic interaction, responsible for transmitting electric and magnetic forces between particles charged under it. The weak interaction, which governs processes like beta decay and neutrino interactions, is mediated by the $W^{+}$, $W^{-}$, and $Z^{0}$ bosons. In contrast, the strong interaction, which binds quarks together within protons, neutrons (and other hadrons) and also holds nuclei together, is mediated by gluons. Contrary to the photon, gluons are charged under the strong force. This causes a phenomenon called self-interaction, leading to extremely complex, non-perturbative structures. This includes the formation of glueballs, hypothetical particles composed entirely of gluons and gluon ladders. Gluon ladders, specifically, will be expanded upon in later section.   Unlike electromagnetism, where the force is mediated by photons that do not interact with each other, the strong interaction is mediated by gluons, which are massless bosons, charged under the strong force, with charge which is referred to as color. Contrary to the electric charge, color charge has 6 unique elements, labeled as red, green, blue and their counterparts (i.e - anti-red). The strong interaction is responsible for binding the nucleons, which are comprised of thee generations of quarks (anti-quarks) and gluons together, such as proton (nominally uud) and neutrons ( nominally udd).

As a result of the symmetry and non-Abelian nature of QCD, there cannot be free particles with color charge. Instead, all observed particles are color-singlets. Contrary to electromagnetism, one of the proposed potential of the strong interaction, is linear with distance at large distances (Cornell potential \cite{brambilla_quark_1999}):
\begin{equation}\label{eq:qcd_potential}
	V(r) = -\frac{4}{3}\frac{\alpha_s}{r} + \sigma r + Const \,,
\end{equation}

where $\alpha_s$ is the strong coupling and $\sigma$ is called the string tension, related to the flux tubes connecting color-charged partons. Due to this potential, as color charged particles drift apart, the potential energy increases, until eventually it is preferable to split the flux tube between the particles and generate new color charged particles in place, as illustrated in Fig. \ref{fig:confinement}.

\begin{figure}[h!] 
	\centering % Centers the image and caption
	\includegraphics[width=0.5\textwidth]{\main/illustrations/color_confinement.png}
	\caption{Illustration of a flux tube breaking, generating new quarks and anti-quarks in the process \cite{greensite_introduction_2020}.}
	\label{fig:confinement}
\end{figure}

\subsection{QCD Jet}\label{sec:qcdjet}
As discussed previously, quarks and gluons are bound by the strong interaction and will generate new color partons when pulled apart. Following that, free quarks and gluons in a highly boosted frame of reference will generate a shower of color partons through a bremsstrahlung cascade,  which will subsequently undergo hadronization (see Fig. \ref{fig:qcd_jet}) and will result in highly collimated clusters of hadrons, leptons and photons \cite{ali_jets_2011}. Additional information regarding the measurable quantity known as jet will be explained in section \ref{sec:reco_jet}.


\begin{figure}[h!] 
	\centering % Centers the image and caption
	\includegraphics[width=0.8\textwidth]{\main/illustrations/Ali_qcd_fragmantation.png}
	\caption{(a) Quark fragmentation to hadrons induced by confinement forces; (b) quark/gluon cascades at high energies in QCD \cite{ali_jets_2011}.}
	\label{fig:qcd_jet}
\end{figure}

The process of hadronization is a topic of active research. Currently research into hadronization uses various models and simulations, as perturbative QCD is not applicable at these energies. 








 
